\documentclass[12pt,a4paper]{article}
% ============================================================
%  Structural Persistence Series — Shared Preamble
% ============================================================
\usepackage{fontspec}
\setmainfont{Hiragino Mincho ProN}
\setsansfont{Hiragino Sans}
\setmonofont{Menlo}

\usepackage{iftex}
\ifXeTeX
  \XeTeXlinebreaklocale "ja"
  \XeTeXlinebreakskip=0pt plus 1pt minus 0.1pt
\fi

\usepackage[margin=22mm]{geometry}
\usepackage{amsmath,amssymb,amsthm}
\usepackage{xcolor}
\usepackage{graphicx}
\usepackage{hyperref}
\usepackage{booktabs}
\usepackage{fancyhdr}
% enumitem not available in this TeX installation

% --- Colours ------------------------------------------------
\definecolor{PaperAccent}{HTML}{1F4E79}
\definecolor{PaperGray}{HTML}{51606F}

\hypersetup{
  colorlinks=true,
  linkcolor=PaperAccent,
  urlcolor=PaperAccent,
  citecolor=PaperAccent,
  pdflang=ja
}
\urlstyle{same}

% --- Typography ---------------------------------------------
\setlength{\parindent}{1.5em}
\setlength{\parskip}{0pt}
\setlength{\emergencystretch}{3em}
\linespread{1.08}
\setlength{\abovedisplayskip}{8pt plus 2pt minus 4pt}
\setlength{\belowdisplayskip}{8pt plus 2pt minus 4pt}
\setlength{\abovedisplayshortskip}{6pt plus 2pt minus 2pt}
\setlength{\belowdisplayshortskip}{6pt plus 2pt minus 2pt}

% --- Section label names ------------------------------------
\renewcommand{\abstractname}{要旨}

% --- Theorem-like environments ------------------------------
\theoremstyle{plain}
\newtheorem{proposition}{命題}[section]
\newtheorem{lemma}[proposition]{補題}
\newtheorem{corollary}[proposition]{系}

\theoremstyle{definition}
\newtheorem{definition}{定義}[section]
\newtheorem{assumption}{条件}[section]

\theoremstyle{remark}
\newtheorem*{remark}{注}

% --- Running header -----------------------------------------
\pagestyle{fancy}
\fancyhf{}
\renewcommand{\headrulewidth}{0.4pt}
\fancyhead[L]{\small\textcolor{PaperGray}{\nouppercase{\leftmark}}}
\fancyhead[R]{\small\textcolor{PaperGray}{\thepage}}
\fancypagestyle{plain}{%
  \fancyhf{}
  \renewcommand{\headrulewidth}{0pt}
  \fancyfoot[C]{\small\textcolor{PaperGray}{\thepage}}
}

% --- Title block macro --------------------------------------
\newcommand{\PaperTitleBlock}[3]{%
  % #1 = title, #2 = subtitle, #3 = date
  \thispagestyle{plain}%
  \begin{center}
  {\LARGE\bfseries #1\par}
  \vspace{0.4em}
  {\normalsize #2\par}
  \vspace{1.0em}
  {\large Akihito Sunagawa\par}
  \vspace{0.3em}
  {\normalsize #3\par}
  \end{center}
  \vspace{1.6em}
}

% Legacy 2-arg version (backward compat)
\newcommand{\PaperTitleBlockSimple}[2]{%
  \PaperTitleBlock{#1}{}{#2}
}


\begin{document}

\PaperTitleBlock
  {有限CSPにおける構造持続の予測力}
  {— Mixed-SAT/NAE-SAT と Exp43c q-coloring による Route A 経験的検証 —}
  {2026年4月12日}

\begin{abstract}
本補論は、構造持続理論の Route A に属する有限制約充足問題において、累積構造消耗 \(L\) が単なる制約数より強い予測情報を持つかを検証する。

単一の制約族、たとえば 3-SAT だけ、または NAE-SAT だけを見るなら、各制約の drift は定数である。したがって

\[
L = m \cdot d
\]

となり、\(L\) は raw constraint count \(m\) の定数倍にすぎない。この設定では、\(L\) が raw count より良いかを問うことは原理的に縮退している。

そこで本補論では、同じインスタンス内に 3-SAT 制約と 3-NAE-SAT 制約を混ぜる Mixed-CSP を用いる。3-SAT の drift は

\[
d_{\mathrm{SAT}}=\log(8/7)\approx 0.1335
\]

であり、3-NAE-SAT の drift は

\[
d_{\mathrm{NAE}}=\log(4/3)\approx 0.2877
\]

である。raw count は両者を同じ 1 制約として数えるが、\(L\) は制約の質を drift-weighted に数える。この差により、構造持続理論の経験的主張を非縮退に検査できる。

2026-04-22 に事前登録済み primary run を実行した。全 12,000 インスタンスが正常に完走し、timeout は 0、malformed encoding は 0 であった。leave-one-mixture-out の予測比較では、primary predictor である \(L+n\) モデルが raw count + \(n\) baseline を大きく上回った。

{\def\LTcaptype{table} % do not increment counter
\begin{longtable}[]{@{}lrrr@{}}
\toprule\noalign{}
model & log loss & Brier & accuracy@0.5 \\
\midrule\noalign{}
\endhead
\bottomrule\noalign{}
\endlastfoot
\texttt{L\_plus\_n} & 0.0970 & 0.0299 & 0.9631 \\
\texttt{cnf\_count\_plus\_n} & 0.1010 & 0.0304 & 0.9631 \\
\texttt{first\_moment} & 0.1489 & 0.0482 & 0.9457 \\
\texttt{raw\_density} & 0.7447 & 0.2524 & 0.6279 \\
\texttt{raw\_plus\_n} & 0.7525 & 0.2539 & 0.6159 \\
\texttt{raw\_count} & 0.7708 & 0.2798 & 0.5139 \\
\end{longtable}
}

事前登録された四つの判定基準はいずれも通過した。

\begin{itemize}
\tightlist
\item
  primary support: \texttt{L\_plus\_n} log loss 0.0970 \textless{} \texttt{raw\_plus\_n} 0.7525
\item
  strong support: \texttt{raw\_plus\_n} に対する相対改善 87.1\%（閾値 10\%）
\item
  theory-pure support: \texttt{first\_moment} 0.1489 \textless{} \texttt{raw\_plus\_n} 0.7525
\item
  encoding guardrail: \texttt{L\_plus\_n} 0.0970 \textless= \texttt{cnf\_count\_plus\_n} 0.1010
\end{itemize}

この結果は、\(L\) が単なる制約数ではなく、制約が解空間を削る質的差を予測情報として運ぶことを示す。これは Bernoulli-CSP universality class に対する経験的支持であり、SAT 単独の結果を finite CSP class 内へ拡張する最初の非縮退な prospective test である。

2026-04-23 には、Exp43c q-coloring primary validation も freeze 済み threshold-local protocol のもとで実行された。4,000 インスタンスが正常に完走し、timeout と malformed encoding は 0 であった。leave-one-q-out の予測比較では、theory-specified predictor \texttt{fm\_plus\_n} が best primary raw baseline を上回った。

{\def\LTcaptype{table} % do not increment counter
\begin{longtable}[]{@{}lr@{}}
\toprule\noalign{}
model & mean held-out log loss \\
\midrule\noalign{}
\endhead
\bottomrule\noalign{}
\endlastfoot
\texttt{fm\_plus\_n} & 0.440189 \\
\texttt{first\_moment} & 0.446814 \\
\texttt{raw\_density} & 0.780910 \\
\texttt{density\_plus\_n\_q} & 2.804019 \\
\texttt{cnf\_count\_plus\_n\_q} & 7.700105 \\
\texttt{raw\_plus\_n\_q} & 8.567224 \\
\end{longtable}
}

Exp43c の結果は、SAT / NAE-SAT mixed CSP だけでなく、SAT 構文の外側に見える graph-coloring family でも、first-moment / drift coordinate が raw / density / encoding-size baselines より out-of-sample に強いことを示す。ただし、これは absolute q-coloring threshold の予測ではなく、frozen threshold-local window 内での feasibility prediction である。
\end{abstract}

\section{目的}\label{ux76eeux7684}

構造持続理論の最小形式では、状態集合の縮小は対数構造消耗として蓄積される。制約充足問題では、状態集合は候補割当ての集合であり、各制約はその集合を一定割合だけ削る。

ランダム 3-SAT では、一様な割当てが単一節を破る確率は \(1/8\) であり、単一節を満たす確率は \(7/8\) である。したがって各節による drift は

\[
d_{\mathrm{SAT}}=-\log(7/8)=\log(8/7).
\]

同様に、3-NAE-SAT では全て真または全て偽の 2 通りが禁止されるため、単一 NAE 制約を満たす確率は \(6/8=3/4\) であり、

\[
d_{\mathrm{NAE}}=-\log(3/4)=\log(4/3).
\]

このとき混合インスタンスの累積構造消耗量は

\[
\begin{aligned}
L \\
&= \\
m_{\mathrm{SAT}}\log(8/7) \\
+ \\
m_{\mathrm{NAE}}\log(4/3).
\end{aligned}
\]

第一モーメント法は、候補割当て数の期待値を

\[
\begin{aligned}
E[\#\mathrm{SAT}] \\
&= \\
2^n e^{-L} \\
&= \\
\exp(n\log 2-L)
\end{aligned}
\]

と与える。したがって

\[
F = n\log 2 - L
\]

は、充足割当て数の第一モーメント対数である。

本補論の問いは、第一モーメントそのものを証明することではない。それは制約族の仕様から従う。問いは、その \(L\) または \(n\log 2-L\) が、raw count や raw density よりも、有限インスタンスの feasibility を out-of-sample に予測するかである。

\section{単一 family での縮退}\label{ux5358ux4e00-family-ux3067ux306eux7e2eux9000}

NAE-SAT 単独で

\[
L=m\log(4/3)
\]

を用いて feasibility を予測しても、これは raw count \(m\) と完全に同値である。3-SAT 単独でも同じである。

したがって、単一 family 内で

\begin{verbatim}
$L normalization > raw count$
\end{verbatim}

を主張することはできない。回帰モデルでは係数が定数倍に変わるだけで、予測性能は同一になる。

非自明な検査には、少なくとも二つの条件が必要である。

第一に、制約ごとの drift が異なること。第二に、評価 endpoint が raw count だけでは区別できない差を含むこと。

Mixed-SAT/NAE-SAT はこの条件を満たす。たとえば \(n=120\)、density \(d=2.5\) では raw count は全 mixture で \(m=300\) に固定される。しかし SAT/NAE の比率を変えると \(L\) は変わる。

{\def\LTcaptype{table} % do not increment counter
\begin{longtable}[]{@{}lrlr@{}}
\toprule\noalign{}
mixture & raw count & qualitative drift & observed SAT rate \\
\midrule\noalign{}
\endhead
\bottomrule\noalign{}
\endlastfoot
pure SAT & 300 & low & 1.000 \\
25\% NAE & 300 & low-medium & 1.000 \\
50\% NAE & 300 & medium & 1.000 \\
75\% NAE & 300 & high & 0.145 \\
pure NAE & 300 & highest & 0.000 \\
\end{longtable}
}

raw count はこれらを全て同一視する。\(L\) は NAE 制約を SAT 制約より重く数えるため、この差を区別できる。

\section{実験設計}\label{ux5b9fux9a13ux8a2dux8a08}

\subsection{インスタンス}\label{ux30a4ux30f3ux30b9ux30bfux30f3ux30b9}

primary grid は以下で構成される。

\begin{itemize}
\tightlist
\item
  \(n\in\{80,120,160\}\)
\item
  density \(d=m/n\in\{2.0,2.5,3.0,3.5\}\)
\item
  mixture:
\end{itemize}

\begin{verbatim}
- pure SAT
- 75% SAT / 25% NAE
- 50% SAT / 50% NAE
- 25% SAT / 75% NAE
- pure NAE
\end{verbatim}

\begin{itemize}
\tightlist
\item
  各 cell 200 インスタンス
\end{itemize}

合計は

\[
3\times 4\times 5\times 200=12000
\]

である。

exact-one-3-SAT は、drift が大きく CNF encoding size との交絡も強いため、初回 primary には含めない。これは保守的な判断である。SAT/NAE の二型混合だけでも \(L\neq \mathrm{constant}\times m\) は成立し、raw count との非縮退比較が可能である。

\subsection{endpoint}\label{endpoint}

primary endpoint は feasibility である。すなわち、インスタンスが SAT か UNSAT かを予測する。

これは重要である。以前の計算コスト補論では、条件つき発見確率 \(P(\mathrm{found}\mid\mathrm{SAT})\) と solver cost を扱った。しかし Route A の第一モーメント構造が直接支配するのは、まず解空間の体積と feasibility であって、特定ソルバーの探索コストではない。

したがって本補論では、solver conflict count ではなく SAT rate を primary に置く。solver cost は別の問いであり、存在論的 \(L\) と発見論的 \(c\) を混同しない。

\subsection{予測モデル}\label{ux4e88ux6e2cux30e2ux30c7ux30eb}

事前登録では、leave-one-mixture-out cross-validation により、以下のモデルを比較した。

\begin{itemize}
\tightlist
\item
  \texttt{raw\_count}: raw semantic constraint count
\item
  \texttt{raw\_density}: \(m/n\)
\item
  \texttt{raw\_plus\_n}: raw count + \(n\)
\item
  \texttt{L\_only}: drift-weighted \(L\)
\item
  \texttt{L\_plus\_n}: \(L+n\)
\item
  \texttt{first\_moment}: \(n\log2-L\)
\item
  \texttt{cnf\_count\_plus\_n}: CNF clause count + \(n\)
\end{itemize}

primary comparison は \texttt{L\_plus\_n\ \textless{}\ raw\_plus\_n} である。\texttt{first\_moment\ \textless{}\ raw\_plus\_n} は theory-pure support とした。\texttt{cnf\_count\_plus\_n} は encoding guardrail であり、\(L\) の勝利が単に CNF 展開サイズの勝利でないかを検査する。

\section{実行整合性}\label{ux5b9fux884cux6574ux5408ux6027}

\subsection{aborted attempt と verifier fix}\label{aborted-attempt-ux3068-verifier-fix}

最初の primary attempt は 563 行で停止した。2 行が \texttt{malformed\_encoding} と判定されたためである。これは CNF encoding の論理バグではなく、verifier 側の false positive であった。

原因は、PySAT / MiniSat が unconstrained variable を返却 model から省略することがある一方で、初期 verifier が全変数 \(1,\dots,n\) の assignment を要求していたことである。SAT/NAE 制約の意味論評価に必要なのは、実際に制約に現れる変数だけである。

修正後 verifier は、semantic constraints に現れる変数集合だけを要求する。修正後、以下の検査を通した。

\begin{itemize}
\tightlist
\item
  archived malformed rows 2 件の regression replay
\item
  default agreement test 1000 instances
\item
  stress agreement test \(n=80,d=2.0\), 500 instances
\item
  stress agreement test \(n=160,d=2.0\), 500 instances
\end{itemize}

いずれも agreement failure は 0 であった。aborted attempt は OSF bundle と git history に forensic archive として残し、official primary analysis からは除外している。

\subsection{official primary}\label{official-primary}

修正後、official primary を別ファイルでゼロから実行した。

\begin{verbatim}
official primary rows: 12000/12000
status: $succeeded = 12000$
timeouts: 0
malformed encodings: 0
SAT feasible rows: 6753
UNSAT rows: 5247
\end{verbatim}

このため、primary analysis は aborted attempt に依存しない。

\section{結果}\label{ux7d50ux679c}

\subsection{primary model comparison}\label{primary-model-comparison}

leave-one-mixture-out の結果は以下である。

{\def\LTcaptype{table} % do not increment counter
\begin{longtable}[]{@{}lrrr@{}}
\toprule\noalign{}
model & log loss & Brier & accuracy@0.5 \\
\midrule\noalign{}
\endhead
\bottomrule\noalign{}
\endlastfoot
\texttt{L\_plus\_n} & 0.0970 & 0.0299 & 0.9631 \\
\texttt{cnf\_count\_plus\_n} & 0.1010 & 0.0304 & 0.9631 \\
\texttt{first\_moment} & 0.1489 & 0.0482 & 0.9457 \\
\texttt{L\_only} & 0.4731 & 0.1601 & 0.7626 \\
\texttt{raw\_density} & 0.7447 & 0.2524 & 0.6279 \\
\texttt{raw\_plus\_n} & 0.7525 & 0.2539 & 0.6159 \\
\texttt{raw\_count} & 0.7708 & 0.2798 & 0.5139 \\
\end{longtable}
}

primary predictor \texttt{L\_plus\_n} は \texttt{raw\_plus\_n} を大きく上回った。

\[
0.0970 < 0.7525.
\]

相対改善は

\[
1-\frac{0.0970}{0.7525}\approx 0.871
\]

であり、事前登録の strong support 閾値 10\% を大幅に超える。

\subsection{first moment}\label{first-moment}

理論的に最も純粋な predictor は

\[
n\log2-L
\]

である。これは係数を固定した第一モーメント対数であり、追加の自由係数を持たない。

この \texttt{first\_moment} も \texttt{raw\_plus\_n} を明確に上回った。

\[
0.1489 < 0.7525.
\]

ただし \texttt{L\_plus\_n} には負ける。これは有限サイズ補正を示唆する。\texttt{first\_moment} は \(n\) と \(L\) の係数を \((\log2,-1)\) に固定するが、\texttt{L\_plus\_n} は有限 \(n\) の偏りを logistic regression の係数で吸収できる。

したがって、\texttt{first\_moment} と \texttt{L\_plus\_n} の差

\[
0.1489-0.0970=0.0519
\]

は、第一モーメント理論からの finite-size deviation を測る経験的余白として読める。

\subsection{encoding guardrail}\label{encoding-guardrail}

\texttt{L\_plus\_n} は \texttt{cnf\_count\_plus\_n} にもわずかに勝った。

\[
0.0970 \le 0.1010.
\]

この margin は小さい。しかし SAT/NAE の二型混合では、それは予想される。CNF encoding では SAT 制約は 1 clause、NAE 制約は 2 clauses で表せる。一方、drift ratio は

\[
\frac{\log(4/3)}{\log(8/7)}\approx 2.15
\]

であり、CNF clause count の比 2.00 と近い。したがって SAT/NAE-only grid では、\(L\) と CNF clause count はほぼ比例する。

この条件で \texttt{L\_plus\_n} が \texttt{cnf\_count\_plus\_n} に勝つことは、方向としては理論と整合するが、margin が狭くなるのは構造的に自然である。

将来、exact-one 制約を stress extension として含めると、この guardrail はより強く分離される。exact-one-3-SAT の drift は

\[
\log(8/3)\approx 0.981
\]

であり、SAT に対して約 7.3 倍である。一方、pairwise CNF encoding の clause count は SAT に対して 4 倍である。ここでは \(L\) と CNF size の比例性がより大きく破れる。

ただし、exact-one は feasibility の飽和と encoding size 交絡が強いため、初回 primary から外した。この判断は、最初の empirical Route A test を清潔に保つための保守的設計である。

\section{理論的位置づけ}\label{ux7406ux8ad6ux7684ux4f4dux7f6eux3065ux3051}

\subsection{SAT 単独から Bernoulli-CSP class へ}\label{sat-ux5358ux72ecux304bux3089-bernoulli-csp-class-ux3078}

ランダム 3-SAT は、構造持続理論にとって最も硬い Route A anchor の一つである。状態集合は \(\{0,1\}^n\)、測度は counting measure、各制約の縮小率は仕様から直接計算できる。

しかし 3-SAT 単独では、\(L=m\log(8/7)\) であり、raw count と縮退する。したがって、3-SAT で \(L\) が有効であることは重要だが、「制約の質を drift-weighted に数えることが raw count より強い」という主張はまだ検査されていない。

Mixed-SAT/NAE-SAT はこの間隙を埋める。raw count を固定したまま、constraint type の drift を変えることで、\(L\) の非自明な予測力を検査する。

今回の Mixed-CSP 結果は、Bernoulli-CSP universality class に対して、Lean での formal layer に対応する empirical counterpart を与える。具体的には \texttt{Survival.BernoulliCSPUniversality}、\texttt{Survival.NAESATChernoffCollapse}、および関連する exposure / Chernoff wrapper 群に対応する経験的検査である。すなわち、固定割当てに対する Bernoulli bad-event exposure として表現できる異なる CSP 制約族が、同じ \(L\)-based coordinate で feasibility を整理される。

Exp43c q-coloring は、この方向を graph-coloring family へ広げる。Lean 側の q-coloring formalization は fixed-coloring edge exposure を扱う。一方、Exp43c の empirical endpoint は full graph q-colorability である。したがって Exp43c は q-coloring threshold theorem ではない。より正確には、fixed-coloring exposure から得られる first-moment / drift coordinate

\[
n\log q - L
\]

が、threshold-local な q-colorability feasibility を raw edge count、density、CNF encoding-size baseline よりよく予測した、という経験的検査である。

\subsection{何を主張し、何を主張しないか}\label{ux4f55ux3092ux4e3bux5f35ux3057ux4f55ux3092ux4e3bux5f35ux3057ux306aux3044ux304b}

本補論が主張するのは次である。

\begin{enumerate}
\def\labelenumi{\arabic{enumi}.}
\tightlist
\item
  SAT/NAE mixed finite CSP では、\(L+n\) が raw count + \(n\) より feasibility を強く予測した。
\item
  \(n\log2-L\) という theory-pure predictor も raw baseline を上回った。
\item
  CNF encoding size baseline との比較でも、\(L+n\) は少なくとも同等以上であった。
\item
  Exp43c q-coloring では、\(n\log q-L\) に対応する \texttt{fm\_plus\_n} が raw / density / CNF-size baselines を leave-one-q-out で上回った。
\item
  したがって、\(L\) または first-moment coordinate は単なる制約数ではなく、制約の質的 drift を含む予測座標として機能する。
\end{enumerate}

本補論が主張しないのは次である。

\begin{enumerate}
\def\labelenumi{\arabic{enumi}.}
\tightlist
\item
  すべての CSP で同じ定数係数が成立する、とは主張しない。
\item
  drift から solver cost の感度指数 \(c\) を一点予測できる、とは主張しない。
\item
  XOR-SAT のように多項式構造が支配する family を primary empirical cost test に使える、とは主張しない。
\item
  Bootstrap percolation や branching process がこの意味で Route A anchor である、とは主張しない。
\item
  独立再現なしに universal law が確立した、とは主張しない。
\end{enumerate}

ここで得られたのは、Bernoulli-CSP class 内での empirical universality support である。より強い universal law of tendency には、formal theorem 4 の昇格と独立再現が必要である。

\subsection{LLM 実験との対応}\label{llm-ux5b9fux9a13ux3068ux306eux5bfeux5fdc}

LLM 側の Exp.40/41/42 は、scope-as-repair が quality-blind baseline より予測力を持つことを示した。Mixed-CSP は、それとは全く異なる hard combinatorial domain で、drift-weighted \(L\) が raw baseline より予測力を持つことを示す。

両者の機構は同じではない。LLM の attribution-as-repair と SAT/NAE の bad-event exposure は別の現象である。

共通しているのは、単純な「量」ではなく、構造の質を区別する座標が out-of-sample の予測力を持つ点である。LLM では scope / attribution の質が効き、Mixed-CSP では constraint drift の質が効く。

この意味で、Mixed-CSP は「例を一つ足した」のではなく、方法論を横断した検査である。すなわち、quality-blind / raw-count baseline に対して、structure-aware coordinate が prospective に勝つかを問う同じ型の検査である。

\section{限界}\label{ux9650ux754c}

第一に、Mixed-CSP 結果は SAT/NAE の二型混合に限定される。Exp43c は q-coloring への拡張を与えたが、これは単独で universal law を確立するものではない。Cardinality-SAT への拡張は自然な次段階だが、本補論の validated claim ではない。

第二に、encoding guardrail の margin は小さい。これは SAT/NAE では CNF clause ratio と drift ratio が近いためである。この点は exact-one stress extension でより強く検査できる。

第三に、対象は feasibility であり、solver cost ではない。solver cost については、既存の計算コスト補論が存在と発見を分けて扱っている。本補論の結果を、任意ソルバーの runtime 予測へ直接外挿してはならない。

第四に、有限サイズ効果が残る。\texttt{first\_moment} が raw baseline を大きく上回る一方で \texttt{L\_plus\_n} に負けることは、理論的第一モーメントだけでは finite \(n\) の補正を完全には吸収できないことを示す。

第五に、公式 primary の前に verifier false positive による aborted attempt が存在する。ただしこれは archive され、原因は unconstrained variable の model omission と特定され、修正後の regression / agreement test を通過した。official primary は修正後に別ファイルでゼロから実行され、aborted attempt は分析から除外されている。

\section{再現性}\label{ux518dux73feux6027}

関連ファイルは以下にある。

{\def\LTcaptype{table} % do not increment counter
\begin{longtable}[]{@{}
  >{\raggedright\arraybackslash}p{(\linewidth - 2\tabcolsep) * \real{0.5000}}
  >{\raggedright\arraybackslash}p{(\linewidth - 2\tabcolsep) * \real{0.5000}}@{}}
\toprule\noalign{}
\begin{minipage}[b]{\linewidth}\raggedright
file
\end{minipage} & \begin{minipage}[b]{\linewidth}\raggedright
role
\end{minipage} \\
\midrule\noalign{}
\endhead
\bottomrule\noalign{}
\endlastfoot
\texttt{analysis/route\_a\_mixed\_csp/mixed\_csp\_preregistration.md} & 事前登録 \\
\texttt{analysis/route\_a\_mixed\_csp/run\_mixed\_csp.py} & 実行 runner \\
\texttt{analysis/route\_a\_mixed\_csp/analyze\_mixed\_csp.py} & primary analysis \\
\texttt{analysis/route\_a\_mixed\_csp/debug\_mixed\_csp\_encoding.py} & verifier regression / agreement diagnostics \\
\texttt{analysis/route\_a\_mixed\_csp/mixed\_csp\_primary\_official\_2026-04-22.jsonl} & official primary raw records \\
\texttt{analysis/route\_a\_mixed\_csp/mixed\_csp\_results.json} & machine-readable results \\
\texttt{analysis/route\_a\_mixed\_csp/mixed\_csp\_results\_summary.md} & human-readable results \\
\end{longtable}
}

Exp43c q-coloring の関連ファイルは以下にある。

{\def\LTcaptype{table} % do not increment counter
\begin{longtable}[]{@{}
  >{\raggedright\arraybackslash}p{(\linewidth - 2\tabcolsep) * \real{0.5000}}
  >{\raggedright\arraybackslash}p{(\linewidth - 2\tabcolsep) * \real{0.5000}}@{}}
\toprule\noalign{}
\begin{minipage}[b]{\linewidth}\raggedright
file
\end{minipage} & \begin{minipage}[b]{\linewidth}\raggedright
role
\end{minipage} \\
\midrule\noalign{}
\endhead
\bottomrule\noalign{}
\endlastfoot
\texttt{analysis/exp43\_qcoloring/exp43c\_threshold\_local\_preregistration\_draft.md} & fresh threshold-local preregistration \\
\texttt{analysis/exp43\_qcoloring/exp43c\_calibration\_closeout.md} & calibration closeout \\
\texttt{analysis/exp43\_qcoloring/exp43c\_freeze\_package.md} & frozen primary package \\
\texttt{analysis/exp43\_qcoloring/exp43c\_primary\_report.md} & primary validation report \\
\texttt{analysis/exp43\_qcoloring/config/exp43c\_primary\_config.json} & primary grid \\
\texttt{analysis/exp43\_qcoloring/src/evaluate\_primary.py} & frozen evaluation script \\
\end{longtable}
}

OSF addendum:

\begin{itemize}
\tightlist
\item
  zip: \url{https://osf.io/download/69e826573b65e7b53bfd8b7e/}
\item
  manifest: \url{https://osf.io/download/69e8265a30357781bafd90d6/}
\end{itemize}

公式 primary の commit は

\begin{verbatim}
1dadb73dd435347977c47be066002460f136ea00
\end{verbatim}

であり、結果導線の commit は

\begin{verbatim}
a5fc8e2
\end{verbatim}

である。

\section{結論}\label{ux7d50ux8ad6}

Mixed-CSP primary は、構造持続理論の Route A empirical test として、事前登録された四つの判定基準を全て通過した。

3-SAT と 3-NAE-SAT は raw count では同じ 1 制約として見えるが、解空間を削る drift は異なる。\(L\) はこの差を加法的に累積し、finite CSP の feasibility を out-of-sample に予測した。

これは、構造持続理論における \(L\) が単なる量ではなく、制約の質を含む座標として機能することの経験的証拠である。SAT 単独の Route A anchor は、Mixed-SAT/NAE-SAT により Bernoulli-CSP class へ一段拡張された。

さらに Exp43c q-coloring により、SAT 構文の外側に見える graph-coloring family でも、first-moment / drift coordinate が raw / density / encoding-size baseline を上回ることが示された。これは Route A width を一段広げる primary evidence である。

ただし、本補論は universal law の最終宣言ではない。正確な位置づけは、SAT、LLM contradiction repair、Mixed-CSP feasibility、Exp43c q-coloring feasibility で、structure-aware coordinate が quality-blind / raw baseline を上回った、という Level 2 universality candidate への強い支持である。

残る課題は、formal theorem 4 による tendency law への昇格、Cardinality-SAT などへのさらなる幅拡張、そして独立再現である。
\end{document}
